\section{E. The limit at the puncture}
\subsection{E.1. The untwisted pinching and its local smoothings}

In this section $S$ is the smooth minimal rational elliptic surface with
Weierstrass equation
\[
 y^2=x^3-3t^3(t-1)x+2t^4(t-1)^2,
\]
$\pi:S\to B$ is its elliptic fibration, $P$ is the section used in the
period dictionary, $T=t_{6P}:S\to S$, and $L=M_S=\OO_S(P-O)$.
We use the convention that $\mathbb P_S(\OO\oplus L)$ contains the total
space of $L$, with fibre coordinate $a$, zero section $s_0$, and infinity
section $s_\infty$. The two normal bundles are respectively $L$ and
$L^{-1}$. We write $R=\mathbb P_S(\OO\oplus L)$.

\begin{proposition}[E.1.1; Proved]\label{st:E.1.1}
The quotient
\[
 Q=R/\{s_0(x)\sim s_\infty(Tx):x\in S\}
\]
has a compact, reduced complex-space structure. Its normalization is
the finite holomorphic map $\nu:R\to Q$. Its singular locus is a copy
of $S$, and its local analytic equation along that locus is $ab=0$,
with two further coordinates on $S$. The induced map $Q\to B$ is
proper. Its sheaf of first-order local smoothing parameters is
\[
 \mathcal T^1_Q=\mathcal Ext^1_{\OO_Q}(\Omega_Q,\OO_Q)
       \simeq L\otimes T^*L^{-1}\simeq\pi^*\OO_B(1)
\]
on the double locus.
\end{proposition}

\begin{proof}
The two sections of $R$ are disjoint closed submanifolds. The
equivalence relation consists of their two graphs and the diagonal;
it is closed in $R\times R$. A quotient of a compact Hausdorff space
by a closed equivalence relation is Hausdorff: for distinct equivalence
classes, compactness gives disjoint saturated neighbourhoods by taking
finite subcovers of the complements of the relation. Thus the displayed
quotient is compact and Hausdorff.

For a coordinate neighbourhood $U\subset S$ and frames of $L|_U$ and
$L^{-1}|_{T(U)}$, the two tubular coordinate charts are $(x,a)$ at
$s_0(U)$ and $(Tx,b)$ at $s_\infty(T(U))$. We give their identification
the ring
\[
 \{(f(x,a),g(x,b)):f(x,0)=g(x,0)\}
   =\C\{x_1,x_2,a,b\}/(ab).
\]
Changes of frame are holomorphic and invertible on each branch and
have the same change of $x$ on their intersection. They therefore glue
these hypersurface charts to the inherited charts off the sections.
This constructs the asserted complex space. The direct sum of the two
regular branch rings is the integral closure: an integral element in
the total ring of fractions is integral on each smooth branch, and
conversely the two branch idempotents are integral. The resulting map
$\nu$ is proper with at most two points in each fibre, and its local
ring extension is finite. The map to $B$ is well defined because
$\pi T=\pi$; its properness follows from compactness of $Q$.

The conormal presentation for the hypersurface $ab=0$ gives
$\mathcal Ext^1(\Omega,\OO)=\OO/(a,b)$. A change of normal frames
$a'=g(x)a$, $b'=h(x)b$ changes the equation and its smoothing parameter
by $gh$. Consequently these local modules glue as
$N_{s_0/R}\otimes T^*N_{s_\infty/R}=L\otimes T^*L^{-1}$.
For this particular surface, Engel's Lemma 2.1 states
$\operatorname{Hom}(T^*L,L)=\pi^*\OO_B(1)$ \cite[Lemma 2.1]{E26}.
Its identification applies here with exactly $T=t_{6P}$ and $L=\OO(P-O)$;
there is no replacement of the translation by its inverse in this
normal-bundle calculation. One can recover the degree in that lemma
directly: the bundle has multidegree zero on every fibre, since $T$
preserves every component; hence it is $\pi^*\OO_B(n)$. Its degree on
$O$ is
\[
 n=(O-P)\cdot 6P+(P-O)\cdot O=(2-2)+(0-(-1))=1.
\]
The section intersections used here are the explicit computations in
the proof of that lemma. This derives the displayed sheaf in the
normal-bundle convention of our pushout.
\end{proof}

\begin{proposition}[E.1.2; Proved]\label{st:E.1.2}
Over $S^\circ=S\setminus S_\infty$, the pushout in E.1.1 is the central
fibre of a holomorphic plumbing with parameter $u$. On every relatively
compact open disc $U\Subset B\setminus\{\infty\}$ this plumbing is
proper and flat over $U\times\Delta_\epsilon$, for a sufficiently small
positive $\epsilon$. Its fibre at $u\ne0$ is, over $U$, Engel's untwisted
quotient by the contracting generator $u\phi$ covering $T^{-1}$.
This assertion includes discs containing $0$ or $1$; it is made before
the logarithmic transformations at those points.
\end{proposition}

\begin{proof}
Let $\phi$ be the section of $L\otimes T^*L^{-1}$ with divisor
$S_\infty$, as in Corollary 2.2 of \cite{E26}. On $S^\circ$ it is
nowhere zero. At the pair $s_0(x),s_\infty(Tx)$ we use the node chart
\[
       a\otimes b=u\phi_x,
       \qquad a\in L_x,\quad b\in L^{-1}_{Tx}.                 \tag{E1}
\]
This equation is invariant under changes of the two line-bundle
frames, so it defines a hypersurface chart over $S^\circ\times\Delta$.
We describe the gluing, including the required quotient comparison.
Choose a Hermitian metric on $L$ over $\pi^{-1}(\overline U)$, and
$0<r<1/4$. On this compact set the norm of $\phi$ has a finite upper
bound. Shrink $\epsilon$ so that $|u|\|\phi_x\|<r^2$.
Retain the core $r<\|a\|<r^{-1}$ of each projective fibre. Attach the
node chart with $\|a\|<2r$ and $\|b\|<2r$, using the core coordinate
$a$ on its first collar and the coordinate $1/b$ over $Tx$ on its
second collar. The metric used for $b$ is the dual metric over $Tx$.
The two attachment collars in a node chart are disjoint because
$\|a\|\|b\|<r^2$. The transition functions and their inverses are
holomorphic on the overlaps.

These charts are Hausdorff after gluing. Indeed, the smaller closed
core $r\leq\|a\|\leq r^{-1}$ and the closed subcharts
$\|a\|,\|b\|\leq 3r/2$ are compact over a compact subset of
$U\times\Delta_\epsilon$ and have precisely the prescribed closed
identifications on their collars. They cover after a harmless
enlargement of the open cores. The finite closed-chart description
separates distinct points, as in E.1.1, and proves properness over
$U\times\Delta_\epsilon$.

For $u\ne0$, solving (E1) gives
\[
              a_{Tx}=1/b=a_x/(u\phi_x).
\]
This is exactly the inverse of the contracting bundle map
$L_{Tx}\to L_x$, $a_{Tx}\mapsto u\phi_xa_{Tx}$. Thus the map from
the annular charts to $L^\times/\langle u\phi\rangle$ is holomorphic
and identifies exactly the orbits of this generator. Every orbit
meets the core or an attached collar by contraction, and two such
representatives are related by the displayed chart identification.
The comparison map is consequently a biholomorphism. At $u=0$,
(E1) is the pinching chart $ab=0$.

The total node chart is smooth because the derivative of its defining
equation with respect to $u$ is $-\phi_x\ne0$. The projection to $u$
is flat: its local rings are torsion free over the convergent
one-variable power-series ring. The joint map to $(t,u)$ is also
flat. On the core it is a product of $\pi$ with the parameter disc;
$\pi$ is flat since a nonconstant map from a smooth connected surface
to a smooth curve has torsion-free local rings over the curve. On a
node chart, $t-t_0$ and $u-u_0$ form a regular sequence in the regular
total local ring: after setting $t=t_0$, the equation (E1) has no
component contained in a level set $u=u_0$. The two elements cut
dimension by two. The local flatness criterion for a regular base
(equivalently its elementary regular-sequence proof) gives flatness
over $(t,u)$. This argument uses smoothness of the surface $S$, not
smoothness of the Kodaira fibre; it therefore applies at every point
of the fibres $IV^*$ and $III$.
\end{proof}

\begin{proposition}[E.1.3; Proved]\label{st:E.1.3}
The first simultaneous local plumbing at the nodal point of $S_\infty$
has equation
\[
                 xy=u zv,\qquad q=zv,                         \tag{E2}
\]
after a holomorphic change of one normal coordinate, where $q$ is a
local coordinate on $B$ at infinity. This hypersurface is flat over
the $u$-disc. For each $u\ne0$, its three-dimensional fibre has an
ordinary double point at $x=y=z=v=0$. The projective modification
defined by the blow-up of $(x,z)$ resolves that point for $u\ne0$.
Its two strict-transform charts are
\[
       x=zr,\quad ry=uv;
       \qquad z=xa,\quad y=uav.                               \tag{E3}
\]
The modified family remains flat over $u$. At $u=0$, the normalization
of its central fibre contains the blow-up of one branch of (E2) along
the smooth curve $x=z=0$; hence this particular modification changes
the normalization of the unmodified pinching.
\end{proposition}

\begin{proof}
The surface $S$ is smooth and its nodal elliptic fibre has local
equation $q=zv$ at its node. The section $\phi$ has a simple zero
along that fibre. Thus $\phi=zv$ times an invertible holomorphic
function. Absorbing the function in one normal coordinate gives (E2).
The convergent power-series ring $\C\{x,y,z,v,u\}$ is factorial,
and $u$ is coprime to $xy-uzv$. Hence $uh\in(xy-uzv)$ implies
$h\in(xy-uzv)$. The quotient has no $u$-torsion. At a nonzero
parameter the identical argument uses $u-u_0$. This proves flatness
over the smooth parameter curve.
For fixed $u\ne0$ the four first derivatives vanish only at the
origin, and the quadratic form $xy-uzv$ is nondegenerate. This proves
the ordinary-double-point assertion by its defining equation.

The two standard affine charts of the blow-up give (E3), after
removing the common factors $z$ and $x$, respectively. For $u\ne0$
the first chart solves for $v$ and the second for $y$, so both are
smooth. The map is an isomorphism away from the origin in the fixed
$u$-fibre: where the ideal fails to be visibly principal, the
hypersurface equation and one of $y,v$ make it principal. The fibre
over the origin is $\mathbb P^1$. The modification is projective
over (E2), since it is a closed subspace of its product with
$\mathbb P^1$. The first strict-transform equation is coprime to
$u-u_0$ in its convergent power-series ring for every $u_0$; the
second chart is smooth with free parameter $u$. Thus flatness over
$u$ persists.

At $u=0$ the first chart has $ry=0$. The branch $r=0$ maps
isomorphically to $x=0$ in the original central fibre. The branch
$y=0$, together with the second chart, is the ordinary blow-up of the
branch $y=0$ in coordinates $(x,z,v)$ along $(x,z)$. This description
follows from the two affine substitutions defining that blow-up and
is an identification of normalizations, not just a birational map.
For $u\ne0$, the base coordinate in the first chart is
$q=zry/u$ and in the second it is $q=xav$. Thus the modification also
exhibits the triple-coordinate strata that the cusp filling requires.
The substitutions and Jacobian are independently checked in
\texttt{code/e\_degeneration.py}.
\end{proof}

\begin{proposition}[E.1.4; Proved]\label{st:E.1.4}
The compact candidate $Q$ of E.1.1 has
\[
 \pi_1(Q)=\Z,\qquad
 (H_i(Q;\Z))_{i=0}^6
   =(\Z,\Z,\Z^8,\Z^7,\Z^{10},0,\Z),\qquad\chi(Q)=12.
\]
These are invariants of the stated candidate; they do not identify it
as a central fibre of the prescribed compact family.
\end{proposition}

\begin{proof}
The two section embeddings are cofibrations: a tubular neighbourhood
retracts onto each section, and the two neighbourhoods are disjoint.
The ordinary pushout is therefore homotopy equivalent to the space
formed from $R$ by attaching $S\times[0,1]$ along
$s_0$ and $s_\infty T$. This comparison follows by collapsing the
intervals in the mapping cylinder; the tubular retractions construct
its homotopy inverse. Both $S$ and $R$ are simply connected. Van
Kampen applied to this mapping-cylinder model gives one free loop
and no additional generator or relation, so $\pi_1(Q)=\Z$.
The same model has the exact sequence
\[
 H_i(S)\xrightarrow{d_i=s_{0*}-s_{\infty*}T_*}H_i(R)
 \longrightarrow H_i(Q)\longrightarrow H_{i-1}(S).             \tag{E4}
\]
We compute its integral maps, including primitivity.

Put $F=[S_t]$, $m=P-O$, and
$v_0=[6P]-[O]-3F\in H_2(S;\Z)$. The intersection computations in
\cite[Lemma 2.1]{E26} give
\[
 F^2=Fv_0=Fm=0,\quad v_0^2=-6,\quad mv_0=-1,\quad m^2=-2,
 \quad OF=1,\quad Ov_0=0,\quad Om=1.                         \tag{E5}
\]
The translation preserves all components of the two reducible
Kodaira fibres. Their root classes span a negative definite space of
rank seven. Its orthogonal complement is generated over $\Q$ by
$F,O,v_0$. Since $T_*O=O+v_0+3F$, preservation of the intersection
form gives, on this complement and hence on the entire lattice,
\[
 T_*a=a+(aF)v_0-\bigl(av_0-3aF\bigr)F.                       \tag{E6}
\]
For clarity, its value on $v_0$ is $v_0+6F$: it has no $O$-component
and has zero intersections with the rank-seven root space, so it is
$\lambda v_0+aF$. The already established identity $T^*m=m-F$ gives
$mT_*v_0=(m-F)v_0=-1$, forcing $\lambda=1$. Orthogonality to
$T_*O$ then forces $a=6$. These
values on $F,O,v_0$ and the root space verify (E6) on a rational basis,
and both sides preserve the actual integral lattice. In particular
$T^*m=m-F$, in agreement with E.1.1.

The rational elliptic surface is the blow-up of $\mathbb P^2$ at nine
points, so its homology in even degrees is $\Z,\Z^{10},\Z$ and
vanishes in odd degrees; its intersection lattice is unimodular.
The projective-bundle theorem follows here by local trivialization
and the fibre hyperplane class: $H_2(R)=H_2(S)\oplus\Z f$ and
$H_4(R)=\Z\oplus H_2(S)$, with $f$ the projective-line fibre.
The hyperplane class that vanishes on $s_0$ restricts to $-m$ on
$s_\infty$. Therefore
\[
 d_2(a)=(1-T_*)a+(mT_*a)f
       =-(aF)v_0+(av_0-3aF)F+(am-aF)f.                       \tag{E7}
\]
The span of $F,v_0,m$ is primitive. To check the only nontrivial
part of this assertion, split off the unimodular plane spanned by
$F,O$; $m-F$ and $v_0$ lie in its negative definite orthogonal
complement. Their Gram determinant is
$\det\left(\begin{smallmatrix}-2&-1\\-1&-6\end{smallmatrix}\right)=11$.
The square of an index of a larger integral lattice would divide
eleven, so the index is one. The span of $F,v_0$ is also primitive:
pairing with $O,m$ gives determinant $-1$. Unimodularity consequently
makes $a\mapsto(aF,av_0,am)$ surjective onto $\Z^3$. The coefficient
matrix in (E7) has determinant $-1$. Thus $d_2$ has primitive image
$\Z v_0\oplus\Z F\oplus\Z f$, rank three, kernel $\Z^7$,
and cokernel $\Z^8$.

The map $d_4$ sends $[S]$ to $s_0[S]-s_\infty[S]$, whose Poincare
dual is $\pi_R^*m$. This is primitive and nonzero since $mO=1$.
Thus $d_4$ is injective with cokernel $\Z^{10}$. The map $d_0$ is
zero. Substitution in (E4) proves every asserted homology group.
The Euler characteristic is independently obtained by stratification:
$\chi(R)-\chi(S)=24-12=12$. The Smith computations and both Euler
computations are recorded in \texttt{code/e\_degeneration.out}.
\end{proof}

\begin{problem}[E.1.5; Open, with the precise remaining construction]\label{st:E.1.5}
The existence of a flat proper filling of the prescribed compact
untwisted family, with central fibre exactly $Q$, remains undecided
here. E.1.1 constructs the compact candidate as an actual complex
space; E.1.2 constructs its smoothing away from infinity, including
the two other special fibres; E.1.3 computes the first unresolved
two-parameter germ and a concrete flat modification of it.

The smallest missing instance is the following analytic gluing
problem at one lift of the node of $S_\infty$. The local model
$xy=uzv$ must be replaced by charts whose restriction to every
$u\ne0$ is the chosen small resolution of the rank-two Mumford
filling, and these charts must be equivariant under both period
generators, with proper quotient and with central normalization $R$.
The explicit blow-up in E.1.3 supplies the required smooth local
general fibre but changes the normalization at $u=0$. We have not
proved that an alternative gluing preserves it, nor proved that every
alternative must change it. Accordingly this computation establishes
an obstruction to the unmodified plumbing argument, and proves
neither falsity nor truth of Conjecture 8.2. A modified compact limit
also requires an equivariant periodic fan and verification of its
quotient near the crossing $(u,q)=(0,0)$; those data have not been
constructed in this work.
\end{problem}

\subsection{E.2. Finite bielliptic limits and the smooth monodromy}

\begin{proposition}[E.2.1; Proved]\label{st:E.2.1}
After the auxiliary base change $w=s^6$, each of the two families of
\emph{reduced} special surfaces $S_j(w)$ has a proper flat
semistable filling by explicit finite free quotients. The central
surface for $j=1$ is a cycle of four ruled surfaces over an elliptic
curve; that for $j=2$ is a cycle of three. They have
\[
\begin{array}{c|ccccc|c}
 &H_0&H_1&H_2&H_3&H_4&\chi\\\hline
 j=1&\Z&\Z\oplus\Z/3&\Z^5&\Z^9&\Z^4&0\\
 j=2&\Z&\Z\oplus\Z/2&\Z^4&\Z^7&\Z^3&0
\end{array}
\]
with integral coefficients. Their fundamental groups are
$\Z^2\rtimes_{R_3}\Z$ and $\Z^2\rtimes_{R_4}\Z$, respectively,
where
\[
 R_3=\begin{pmatrix}0&-1\\1&-1\end{pmatrix},\qquad
 R_4=\begin{pmatrix}0&-1\\1&0\end{pmatrix}.
\]
\end{proposition}

\begin{proof}
We retain the actual product presentations, which also specify the
comparison with the surfaces in the question. Let
$\rho=e^{\pi i/3}$,
\[
 E_j=\C/(\Z+\kappa_j\Z),\quad
 F_1=\C/(\Z(\rho-2)+\Z(-3)),\quad
 F_2=\C/(\Z(i-1)+\Z(-2)),
\]
where $\kappa_1=\beta_c(z_1)+2\mu(z_1)$ and
$\kappa_2=\beta_c(z_2)+3\mu(z_2)$ have negative imaginary parts.
The holomorphic maps to the central tori are induced by
$L_1(\xi,\eta)=(\eta,\xi-\eta/3)$ and
$L_2(\xi,\eta)=(\eta,\xi-\eta/2)$. Their kernels are generated by
\[
 k_1=(2/3,\rho),\qquad k_2=(1/2,i).
\]
The residual affine generators lift to
\[
 h_1(\xi,\eta)=(\xi+\kappa_1/3,(\rho-1)\eta),\qquad
 h_2(\xi,\eta)=(\xi-\kappa_2/4,-i\eta).                     \tag{E8}
\]
These formulas are verified by substituting the four product periods
in $L_j$; the lattice indices are three and two. Moreover
$(\rho-1)\rho-\rho$ and $(-i)i-i$ are periods of the displayed
$F_j$, so the generators commute. They give free groups
$G_1=C_3\times C_3$ and $G_2=C_2\times C_4$. Indeed, equality of
the first coordinates of $k_j^a h_j^b(\xi,\eta)$ and $(\xi,\eta)$
forces $b=0$ from the coefficient of the nonreal period $\kappa_j$,
and then $a=0$ from the real period. Thus the composite quotient
maps are exactly the degree-nine and degree-eight covers of $S_j$,
with the given signs, rather than abstract surface isogenies.

Use the multiplicative coordinate $v=e^{-2\pi i\xi}$ on $E_j$.
Its Tate parameter is
\[
       q_j=e^{-2\pi i\kappa_j}=A_jw^2=A_js^{12},
       \qquad A_j=e^{-2\pi i\kappa_j^{\mathrm{part}}}\ne0.
\]
Here $A_j$ denotes a scalar only within this proof; it is not a
monodromy matrix. Choose specifically
$A_j^{1/12}=\exp(-2\pi i\kappa_j^{\mathrm{part}}/12)$ and put
$r_j=A_j^{1/12}s$. This choice preserves the fractional period
translations in (E8) exactly. The regular Tate model for $q_j=r_j^{12}$
has a central twelve-gon. We recall its charts to make the extension
explicit. An infinite chain has node charts
$x_i y_i=r_j$; on a branch overlap one uses
$x_{i+1}=y_i^{-1}$ and $y_{i+1}=r_jy_i$. The coordinate
$r_j^i x_i$ on a general fibre is unchanged on overlaps.
The shift $i\mapsto i+12$ acts on that coordinate by $r_j^{12}$.
Quotienting a tubular neighbourhood of the chain by this shift gives
the proper twelve-gon family, and on $r_j\ne0$ it is
$\C^*/q_j^\Z$. Compactness follows from twelve projective components
and a compact annular fundamental domain on a general fibre;
the node charts give flatness and smoothness of its total space.

Multiplication by $q_1^{1/3}=r_1^4$ shifts the component index by four.
Multiplication by $q_2^{-1/4}=r_2^{-3}$ shifts it by minus three.
Hence the maps $h_j$ in (E8) extend, together with their rotations
on $F_j$. The maps $k_j$ multiply $v$ by roots of unity and translate
$F_j$, and extend to every projective component. Their relations
hold on the dense set $s\ne0$ and consequently on all charts.
The action is free on the central fibre as well. A nontrivial power
of $h_j$ moves every component and every node. A nontrivial power
of $k_j$ translates $F_j$ by a nonzero torsion point, including
above the nodes where the multiplicative action alone fixes a point.
The finite free quotient is therefore a smooth total space, proper
over $s$, with a reduced normal-crossing central surface. Flatness
descends along the finite etale quotient. The component orbits have
cardinality three and four, leaving four and three components.
Quotienting an individual $\mathbb P^1\times F_j$ by $k_j$ is a
ruled surface over $F_j/\langle\rho\rangle$ or
$F_j/\langle i\rangle$, respectively. Its two distinguished
sections have the same integral homology class: the corresponding
line bundle has finite order and degree zero on the elliptic base.

For the topology, unwrap the component cycle into an infinite chain
and then pass to the universal cover $\C$ of its elliptic base.
The resulting union of simply connected ruled pieces, with
simply connected intersections and a tree as incidence graph, is
simply connected by van Kampen. Its deck group is generated by
the enlarged lattice of translations on $F_j$ and one step round
the component cycle. The latter conjugates the lattice by the
order-three or order-four rotation. This gives the asserted
semidirect products; inverses of the displayed rotation matrices
give isomorphic groups by reversing the cycle generator.

For homology, the normalization has $n=4$ or $3$ components, each
with ranks $(1,2,2,2,1)$, and $n$ elliptic double curves, each with
ranks $(1,2,1)$. There are no triple points. The mapping-cylinder
sequence for the normalization and the double curves applies
integrally. In degree one the incidence matrix is a twisted cycle:
its cokernel is $\operatorname{coker}(R_j-I)$ and its kernel is
zero, because $R_j$ has no eigenvalue one. The degree-zero incidence
matrix contributes an additional $\Z$ to $H_1$. In degree two the
two section classes on a component agree, so the incidence matrix
is the ordinary cycle matrix on their classes. It has kernel
$\Z$ and primitive image of rank $n-1$. Thus $H_2=\Z^{n+1}$,
$H_3=\Z^{2n+1}$, and $H_4=\Z^n$. Finally
$\operatorname{coker}(R_3-I)=\Z/3$ and
$\operatorname{coker}(R_4-I)=\Z/2$. The exact Smith forms of both
full cycle matrices and the independent $2\times2$ determinants
are in \texttt{code/e\_degeneration.out}.
\end{proof}

\begin{proposition}[E.2.2; Proved]\label{st:E.2.2}
For the retained $S^6$ family over the punctured $w$-disc, the
monodromy acts as the identity on $H^*(S^6;\Z)$. The remaining
smooth monodromy is represented explicitly as follows. Set
$c_0=-i(M+1)$ and $c(r)=c_0-2r$, $0\leq r\leq1$. Let
$H:X_{c_0}\to X_{c_0-2}$ be parallel transport in the proper smooth
family over this interval. With the coherent isomorphism from A.3,
\[
       \Phi=F_{2,c_0-2}\circ H:X_{c_0}\longrightarrow X_{c_0}  \tag{E9}
\]
represents the monodromy of the positively oriented $w$-circle.
Its isotopy class is independent of the smooth connection. The
associated oriented homotopy sphere
$\Sigma_\Phi=D^7\cup_\Phi(-D^7)$ reduces the mapping-class
calculation to the explicit residue
\[
  [\Phi]=28\mu(\Sigma_\Phi)
    =\frac{\langle p_1(W)^2,[W,\partial W]\rangle
                 -4\sigma(W)}{32}\pmod {28},                 \tag{E10}
\]
where $W$ is a spin eight-manifold with boundary $\Sigma_\Phi$ and
the Pontryagin class is lifted to relative rational cohomology.
The identification of $\Z/28$ is oriented by a sphere with
$\mu=1/28$.
\end{proposition}

\begin{proof}
The proper submersion has a smooth horizontal distribution; integration
along a compact path exists because the fibres are compact. The
complex orientation varies continuously. Its monodromy therefore
preserves the orientation class in $H^6$ and fixes the constant
class in $H^0$. These are the only nonzero integral cohomology
groups, proving the first assertion with integral coefficients.

Since $e^{-\pi i(c_0-2r)}=e^{-\pi ic_0}e^{2\pi ir}$, the path
$c(r)$ is the lift of the positive circle. The descent gluing at
its endpoint is exactly $F_{2,c_0-2}$, proving (E9). A path of
horizontal distributions gives a smooth path of the resulting
monodromy diffeomorphisms, so their isotopy classes agree. Changing
the coherent descent maps by a member of the connected automorphism
group also leaves this class unchanged. The diffeomorphism
$X_{c_0}\simeq S^6$ is chosen with the complex orientation.

The glued manifold $\Sigma_\Phi$ is simply connected by van Kampen
and has the integral homology of $S^7$ by Mayer--Vietoris; the
Hurewicz and Whitehead arguments give a homotopy sphere. The
gluing comparison $\pi_0\operatorname{Diff}^+(S^6)=\Theta_7$
sends $\Phi$ to this construction. Here we use Cerf's
Corollary~2 in dimension six \cite[p.~76]{Cerf71}: isotopy classes
are the group of sphere gluing maps modulo extension over a disc.
The extension equivalence is exactly the equivalence of their
two-disc gluings. In dimension seven, the simply connected
h-cobordism comparison identifies these classes with oriented
homotopy spheres; the comparison with oriented diffeomorphism
classes is Smale's Theorem~J \cite[pp.~393--394]{Smale61}.
Thus the map being evaluated is the actual boundary gluing of (E9).
Eells--Kuiper's
seven-dimensional invariant is
$(p_1^2-4\sigma)/896$ modulo one and distinguishes the twenty-eight
oriented homotopy-seven-sphere classes
\cite[Section 7, pp.~100--103]{EK62}. Multiplying that formula by
twenty-eight gives (E10). Its hypotheses hold here: a homotopy
seven-sphere is spin and has vanishing $H^3(-;\Q)$ and
$H^4(-;\Q)$, so the relative lift is unique. The same theorem
and its cobordism construction supply a spin coboundary. Thus
(E10) is an invariant of the specific gluing (E9), rather than an
inference from its action on cohomology.
\end{proof}

\begin{problem}[E.2.3; Open, with local fillings and an invariant reduction]\label{st:E.2.3}
We have not assembled a flat proper filling of the compact $S^6$
family over $w=0$. In particular no global central fibre is named
whose Euler number, fundamental group, or integral homology could
be asserted. E.2.1 constructs and computes the two reduced-surface
degenerations after the stated additional base change; it does not
construct a neighbourhood of the nonreduced fibres $3S_1$ and
$4S_2$ in a compact limit.

The smallest missing finite-place instance is the extension of the
order-three affine deck action over a two-parameter plumbing in
$(t,s)$, with prescribed normal character and central multiplicity
three, followed by descent from $s$ to $w=s^6$. The order-four
analogue must retain the normal character of order four. At infinity
the unresolved periodic gluing is the one specified in E.1.5, now
compatible with these finite-place quotients. The computed central
surface homology groups cannot be substituted for the homology of
that unconstructed threefold.

For the mapping-class part, E.2.2 gives the exact clutching map and
its characteristic-number reduction. No Pontryagin-number or
signature calculation for a coboundary of this particular
$\Sigma_\Phi$ has been completed. The smallest undecided numeric
instance is the residue (E10) for the single loop based at
$c_0=-i(M+1)$; all other basepoints in the connected parameter disc
give conjugate monodromy, and $\Theta_7$ is abelian. The identity
on cohomology proved above supplies no value for that residue.
\end{problem}
